\chapter{Implementation: Data Pipelines and the CycleGAN Baseline}
\label{ch:implementation}

This chapter describes what has actually been built so far towards the Stage~2 baseline set out in Chapter~3: the development infrastructure, the data acquisition and preprocessing pipelines for both domains, a data-quality bug discovered in the initial HiRISE corpus and the pipeline rebuilt to fix it, the CycleGAN model and training infrastructure used to establish the baseline evaluated in Chapter~5, and, once that baseline's evaluation surfaced a deeper corpus-diversity problem than the resolution bug alone explained, a geometry-mediated redesign comprising a real-DTM ground-view renderer, a single-image DTM estimator, and a separate texture-synthesis stage (Section~\ref{sec:geometry-impl}).

\section{Development Infrastructure}

Model training requires a CUDA-capable GPU that the primary development machine does not have, so a dual-machine workflow was set up: development and testing happen locally, while training runs on a departmental A4000 GPU workstation (\texttt{otter70.eps.surrey.ac.uk}). Code moves between the two machines via a private GitHub repository over SSH; the training machine holds a read-only deploy key so it can pull code but never push, matching its role as a training box rather than a development environment. Large data --- raw downloads, processed patches, CTX tiles, and model checkpoints --- is deliberately excluded from version control, since GitHub's free tier cannot sensibly hold multi-gigabyte binaries; each machine instead regenerates its own copy of the data via the download/preprocessing scripts described below. Small, genuinely useful provenance records --- the training log and FID evaluation results --- are the tracked exception, alongside periodic training-progress sample images, so that training runs remain auditable from the repository history alone without needing to re-run them.

\section{Initial Data Pipeline and CycleGAN Implementation}
\label{sec:initial-pipeline}

The initial HiRISE and rover corpora were built by downloading HiRISE browse-quality JPEGs and rover NAVCAM raw images from the NASA Planetary Data System, then preprocessing both into normalised $256{\times}256$ patches: a 1st/99th-percentile intensity stretch, an entropy-based filter (discarding low-texture patches below a fixed entropy threshold, since blank sky or featureless terrain patches contribute nothing useful to training) via \texttt{preprocess.py}, and an 80/10/10 train/val/test split. This produced an initial corpus of 37,054/4,631/4,633 HiRISE patches and 20,417/2,552/2,553 rover patches.

The CycleGAN model (\texttt{cyclegan.py}, \texttt{train\_cyclegan.py}) follows the standard architecture of Zhu \textit{et al.}: two ResNet-9 generators (11.36M parameters each, encoder--9~residual-blocks--decoder with instance normalisation and reflection padding) and two PatchGAN-70 discriminators (2.76M parameters each, outputting a $30{\times}30$ patch-level realism map), trained with LSGAN loss, cycle-consistency weight $\lambda_{\mathrm{cyc}}=10$, identity-loss weight $\lambda_{\mathrm{idt}}=5$, Adam ($\mathrm{lr}=2{\times}10^{-4}$, $\beta_1=0.5$), a 50-image replay buffer per discriminator to reduce oscillation, and mixed-precision training. A quantitative evaluation pipeline was also built alongside the model: \texttt{generate\_translations.py} runs a trained generator over the HiRISE test set, \texttt{prepare\_real\_rgb.py} converts the real rover test set to a matching RGB format, and \texttt{compute\_fid.py} computes the Fr\'{e}chet Inception Distance between the two, using \texttt{clean-fid}'s InceptionV3 feature extractor with a locally reimplemented Fr\'{e}chet-distance formula (the installed \texttt{clean-fid} version's own implementation is incompatible with the environment's current \texttt{scipy}).

\section{A Data-Quality Bug and the Real-Resolution Pipeline Rebuild}
\label{sec:fullres-pipeline}

An early training run (25 epochs, on the initial corpus) completed and was evaluated at FID~171.8. Before proceeding further, an audit of the actual HiRISE imagery being used revealed a data-quality bug: the corpus was built from HiRISE \emph{browse} JPEGs (\texttt{PDS/EXTRAS/RDR}, effectively $\sim$3\,m/pixel after JPEG compression and downsampling), not the true 25\,cm/pixel radiometric data product (RDR, \texttt{PDS/RDR}) the methodology specifies. Continuing to train and evaluate on browse-quality imagery would silently misrepresent the resolution the dissertation's Stage~2 baseline is meant to operate at, so the corpus was rebuilt from genuine full-resolution products rather than continuing to iterate on the bugged data.

The corrected pipeline (\texttt{hirise\_fullres.py}) differs from the original browse-image downloader in several respects, driven by two hard constraints: the training machine's home directory is NFS-backed and quota-limited (a few gigabytes free), and each real RDR product is a single $\sim$568\,MB 16-bit JP2 file, so at most one can be on disk at a time.

\begin{itemize}
  \item \textbf{Real RDR URL construction.} Only \texttt{\_RED.JP2} entries (single-band, full-resolution, non-colour products) are used, filtered from the HiRISE footprint index; this mirrors the existing browse-image filter but points at the true RDR product tree rather than the browse-image tree.
  \item \textbf{Verified download-then-delete lifecycle.} Each observation's JP2 is downloaded, its size checked exactly against the HTTP \texttt{Content-Length} header (a truncated download is deleted and logged as failed rather than silently processed), used to extract patches, and then deleted --- via a \texttt{try/finally} that runs on both success and failure --- before the next observation begins, so disk usage never exceeds one source file at a time.
  \item \textbf{Randomised, entropy-filtered patch sampling.} Patch positions are sampled randomly (with a fixed seed for reproducibility) across each image's full extent rather than filled sequentially from one corner, then filtered by the same entropy threshold as the original pipeline, up to a target of 5,000 qualifying patches per observation (accepting fewer, rather than lowering the entropy threshold, if an observation does not yield enough).
  \item \textbf{Shared split logic.} The train/val/test splitting logic already used by the browse-image pipeline was extracted into a reusable function (\texttt{split\_and\_move}) so both pipelines split patches identically rather than duplicating that logic.
  \item \textbf{CLI orchestration and corpus replacement.} A command-line entry point selects real-RED-JP2 observations within a target region (Oxia Planum) from the HiRISE footprint index, processes them one at a time, and --- once all observations are processed --- replaces the existing browse-resolution HiRISE corpus with the new one, writing a manifest recording each observation's status and patch yield.
\end{itemize}

Running this pipeline against ten targeted Oxia Planum observations, seven yielded usable RED JP2 products (the remaining three were not present as real RED JP2 entries in the footprint index for the requested region), all seven processed successfully with no download or read failures, producing 30,229 patches (24,183/3,022/3,024 train/val/test) against a target of $\sim$50,000 --- a smaller corpus than targeted, but one built entirely from genuine 25\,cm/pixel radiometric data rather than compressed, resampled browse imagery. This corpus, not the original browse-resolution one, is what the training run reported in Chapter~5 uses.

\section{The Geometry-Mediated Redesign}
\label{sec:geometry-impl}

Chapter~5 shows that the corrected, real-resolution corpus of Section~\ref{sec:fullres-pipeline} produced a CycleGAN run with \emph{worse} FID than the buggy browse-resolution one, and traces this to a discriminator shortcut. A follow-up audit of that corpus went a layer deeper: of the ten Oxia Planum observations the full-resolution pipeline targeted, only seven yielded a usable RED JP2 product, collapsing the training corpus to a single region rather than the four regions (299 observations, 37{,}054 patches) the original browse-resolution corpus had covered. A parallel review of the super-resolution, image-to-image translation, and cross-view synthesis literature, conducted in response, surfaced a pattern absent from the original three-stage design: published methods that translate nadir/orbital imagery into ground-level/perspective imagery do not attempt pure pixel-space translation --- they insert an explicit geometric intermediate (a height field or DTM) between the nadir input and the perspective output, and synthesise texture only in the already-correctly-projected view. This section describes the three new components built in response: a real-DTM ground-view renderer, a geometry-mediated corpus assembler, and a single-image DTM estimator, together with the texture-synthesis stage and end-to-end orchestration used to chain them. Chapter~5 reports the results of training and evaluating each.

\subsection{Real HiRISE Stereo DTM Sourcing and the Ground-View Renderer}
\label{sec:renderer-impl}

\texttt{fetch\_hirise\_dtm.py} downloads a HiRISE stereo Digital Terrain Model product (a co-registered height raster, typically 1--2\,m/pixel, plus its source orthoimages) for a given PDS product ID, reusing the same download-then-delete lifecycle as the full-resolution HiRISE pipeline (Section~\ref{sec:fullres-pipeline}) since a single DTM+ortho triplet is itself several hundred megabytes. A PDS ODE coverage query restricted to the Oxia Planum landing ellipse found 27 candidate stereo DTM products in a broader search bounding box, of which 8 fall genuinely within the landing region rather than neighbouring terrain (Mawrth Vallis, Arabia Terra); only these 8 were used, accepting a smaller but geographically-correct corpus over a larger, content-diluted one.

\texttt{dtm\_arrays.py} loads a DTM and its co-registered ortho into aligned NumPy arrays via \texttt{rasterio}, resampling if the two products' resolutions differ and masking the DTM's nodata sentinel to \texttt{NaN} rather than treating it as a real height value. \texttt{render\_ground\_view.py} then projects this height field into a simulated ground-level camera view via height-field ray-marching --- no mesh or GPU rasteriser, since the training machine has no admin rights to install one: for each output column, rays are marched outward from a camera position at a fixed height above the local terrain (1.2\,m, matching Curiosity NAVCAM) and pitch, the nearest un-occluded terrain intersection is found, and the source ortho's pixel value at that world position is drawn as an albedo/shading proxy. This is deliberately the only geometric step in the pipeline: it guarantees correct large-scale shape and viewpoint by construction, rather than asking an adversarial network to infer both from a raw nadir crop, which Chapter~5 (Section~\ref{sec:track1}) shows a 2D convolutional generator cannot do regardless of loss tuning.

\subsection{Geometry-Mediated Corpus Assembly}
\label{sec:geocorpus-impl}

\texttt{assemble\_geometry\_corpus.py} orchestrates dense random-camera-position rendering across the 8 landing-region DTMs, one product on disk at a time. An entropy quality filter carried over from the original preprocessing pipeline initially rejected nearly all candidates (12/800 on a single-product validation run) because it measured whole-frame entropy, and low-relief Oxia Planum terrain commonly fills over half the frame with unrendered ``sky'' (any ray that never reaches terrain) --- diluting the measurement regardless of how textured the actually-rendered ground was. A revised \texttt{ground\_entropy} metric, which excludes sky-value pixels before measuring entropy, fixed this: the same product went from 12/800 to 200/200 (the per-product cap) qualifying crops. The full run across all 8 DTMs produced 1{,}513 images (1{,}210/151/152 train/val/test), with sky-pixel fraction averaging 54.5\% per image, confirming the filter targets genuine ground texture rather than passing degenerate frames. This corpus, not the raw HiRISE crop, is the training input for the Stage~2 CycleGAN retrain reported as Track~2 in Chapter~5.

\subsection{Single-Image DTM Estimation}
\label{sec:dtmestimator-impl}

The renderer above needs a real stereo DTM as input, but stereo DTM coverage only exists for a handful of products per site. \texttt{train\_dtm\_estimator.py} forks the Stage~2 texture-synthesis ControlNet (Section~\ref{sec:controlnet-impl}) into a ControlNet+LoRA model that instead takes a single HiRISE ortho patch and predicts a per-pixel height map, trained on a real, held-out-by-product Gale Crater DTM/ortho corpus (Gale rather than Oxia Planum, since it is where real rover-pose target photography also exists, letting the same product support both this estimator's evaluation and the end-to-end chain in Section~\ref{sec:e2e-impl}). Height targets are encoded per-patch to a normalised unit range (\texttt{height\_encoding.py}) rather than a single global scale, since absolute elevation is not a meaningful quantity for a single 256$\times$256 crop.

\subsection{Perspective Texture Synthesis and End-to-End Orchestration}
\label{sec:controlnet-impl}
\label{sec:e2e-impl}

A separate ControlNet checkpoint (referred to as Track~3 throughout Chapter~5, to distinguish it from the two CycleGAN tracks) is fine-tuned on a paired corpus of real DTM-derived crude renders and their real rover-photo counterparts --- 27 pairs, assembled by pairing a real rover pose's photograph with a render of the real DTM from that same recovered pose, the only source of genuinely paired condition/target imagery in the whole project. Conditioned on a crude geometric render (real or estimator-produced), it synthesises the fine surface texture the renderer's flat albedo draping cannot represent, playing the same photorealism role the original design's UNSB physics losses were meant to constrain (Chapter~3, Section~\ref{sec:formulation}) but via paired supervision on a small real corpus rather than an unpaired physics-informed objective.

\texttt{run\_end\_to\_end\_pipeline.py} chains all three pieces for the first time: a HiRISE ortho patch through the DTM estimator, the renderer (at a fixed synthetic camera pose, run twice --- once against the real stereo DTM, once against the estimator's own output), and the Track~3 ControlNet, so that the question Chapter~5 (Section~\ref{sec:e2e-results}) asks --- does the DTM estimator's height error visibly propagate to the final image, relative to using the real DTM --- has an actual image answer rather than an inference from three separately-evaluated pieces. Because the real DTM/ortho fetch (\textasciitilde940\,MB) and the GPU inference it feeds do not need to happen on the same machine, \texttt{prepare\_e2e\_crop.py} splits the fetch out as a separate, disk-heavy, GPU-free step that ships only the resulting \textasciitilde320\,KB cropped array to wherever inference actually runs.

\section{Summary}

The infrastructure, data pipelines, model, and evaluation tooling needed to train and assess a Stage~2 baseline are now in place and have been exercised end-to-end: the original CycleGAN pipeline (Sections~\ref{sec:initial-pipeline}--\ref{sec:fullres-pipeline}), and, built in response to that pipeline's diagnosed failure, a geometry-mediated redesign comprising a real-DTM renderer, a geometry-mediated corpus assembler, a single-image DTM estimator, and a texture-synthesis ControlNet chained end-to-end (Section~\ref{sec:geometry-impl}). Chapter~5 reports the results of training and evaluating each component, including recovering from two further data-quality findings discovered partway through rather than continuing to build on them.
